\subsection*{Висновки до першого розділу}
\addcontentsline{toc}{subsection}{Висновки до першого розділу}

Теоретичний аналіз проблеми формування інформаційно-комунікативної компетентності (ІКК) майбутніх дизайнерів дозволяє сформулювати такі висновки.

1. На основі аналізу наукових джерел уточнено поняттєвий апарат дослідження: ІКК визначено як інтегративну характеристику особистості дизайнера, що відображає його здатність ефективно здійснювати пошук, аналіз, оцінювання та використання професійної інформації (у т.\,ч. згенерованої ШІ), комунікувати за допомогою цифрових засобів, застосовувати сучасні цифрові та ШІ-інструменти дизайну, а також рефлексувати щодо власного рівня ІКК та етичних аспектів використання технологій. Обґрунтовано 4-компонентну структуру ІКК, специфічну для дизайнерів: інформаційно-аналітичний, комунікативний, технологічний та рефлексивний компоненти. Проведено розмежування ІКК із суміжними поняттями (ІКТ-компетентність, цифрова компетентність, медіакомпетентність).

2. Оглядове дослідження 156 наукових праць з проблеми генеративного ШІ в дизайн-освіті (Web of Science, 2017--2026) виявило 6,4-кратне зростання публікацій після появи ChatGPT (86,5\% досліджень опубліковано після 2022~р.). Встановлено домінування тематичних кластерів: креативність (35,9\%), оцінювання (27,6\%), дизайн навчальних програм (22,4\%). Критичною є теоретична прогалина: лише 26,3\% досліджень використовували теоретичні рамки. Зафіксовано парадигмальний зсув від <<дизайнера-виробника>> до <<дизайнера-куратора>>, що безпосередньо впливає на зміст усіх компонентів ІКК.

3. Аналіз стандарту вищої освіти В2 Дизайн у контексті Європейської рамки кваліфікацій виявив структурну відповідність за вимірами знань, умінь та автономності, проте з суттєвими прогалинами у компетентностях, пов'язаних з критичною рефлексією, міждисциплінарною співпрацею, сталим дизайном та, найважливіше, цифровими/ШІ-навичками. Студійна (проєктна) педагогіка на основі рефлексивної практики визначена як найбільш теоретично обґрунтований підхід до формування ІКК дизайнерів.

4. Розроблено інтегративну теоретичну рамку, що поєднує модель TPACK (адаптовану для дизайну: TK=ШІ/цифрові інструменти, PK=студійні методи, CK=предметна область дизайну), студійну педагогіку (рефлексивна практика Шона) та компетентнісний підхід (стандарт В2). Ця рамка є наскрізною для дослідження: вона спрямовує аналіз даних у розділі~\ref{sec:2} та обґрунтовує модель трансформації навчальних програм у розділі~\ref{sec:3}.

Виявлені теоретичні засади обґрунтовують необхідність комплексного аналізу поточного стану формування ІКК в українських програмах дизайну на основі даних, що є предметом розділу~\ref{sec:2}.
